\section{Differential equation and local solution bases}\label{sec:ode}

Put $\thetaop=z\frac{\dd}{\dd z}$.  The differential equation listed on
Wikipedia is
\begin{equation}\label{eq:meijer-ode}
 \left[(-1)^{p-m-n}z\prod_{j=1}^{p}(\thetaop-a_j+1)
       -\prod_{j=1}^{q}(\thetaop-b_j)\right]G(z)=0.
\end{equation}

\subsection{The kernel shift identity}

The proof is a one-line Gamma recurrence once the signs are tracked correctly.
From $\Gamma(w+1)=w\Gamma(w)$,
\begin{equation}\label{eq:phi-shift}
 \frac{\Phi(s)}{\Phi(s-1)}
 =(-1)^{p-m-n}
   \frac{\prod_{j=1}^{p}(s-a_j)}
        {\prod_{j=1}^{q}(s-b_j)}.
\end{equation}
The exponent can also be written $p+m-n$; the two differ by the even number
$2m$.

\begin{theorem}[Meijer differential equation]\label{thm:ode}
Every certified presentation of $G_{p,q}^{m,n}$ satisfies
\eqref{eq:meijer-ode} in the quantitative derivative sense of
\cref{thm:analyticity}.
\end{theorem}
\begin{proof}
On a finite vertical contour, $\thetaop$ multiplies $z^s$ by $s$.  Hence the
second product in \eqref{eq:meijer-ode} has integrand
$\prod_j(s-b_j)\Phi(s)z^s$.  In the first product, shift $s$ to $s-1$ after the
extra factor $z$; its integrand becomes
$(-1)^{p-m-n}\prod_j(s-a_j)\Phi(s-1)z^s$, equal by
\eqref{eq:phi-shift}.  The two finite contours differ only by horizontal
connector segments and residues crossed by the shift.  Choosing a common
admissible strip, or performing the argument on the residue normal form, makes
these corrections cancel.  Tail and derivative errors are allocated below
$2^{-k-2}$.  In the nonconfluent series presentation, the same result follows
more directly: inserting $z^{b_h+k}$ turns \eqref{eq:meijer-ode} into the
coefficient recurrence obtained from \eqref{eq:hypergeom}.  For confluent
power-log terms, $\thetaop$ acts on the finite polynomial in $\Log z$, and the
Laurent-jet identity gives the same cancellation.
\end{proof}

\subsection{Order and singularities}

Expanding the products in $\thetaop$ shows that the equation has order
$\max(p,q)$ except at a degenerate leading coefficient.  At $z=0$ it is a
regular singular equation when $p\le q$ in the standard local direction; at
$z=\infty$ the inverted equation is analogous.  When $p=q$, the leading
coefficient vanishes at
\begin{equation}\label{eq:balanced-singularity}
 z=(-1)^{p-m-n},
\end{equation}
which explains the additional singular point mentioned on the source page.
These are finite polynomial calculations in the coefficients of
\eqref{eq:meijer-ode}.

\subsection{Fundamental systems and the missing hypothesis}

For $p\le q$, place $b_h$ first in the bottom list and consider
\begin{equation}\label{eq:right-basis}
 Y_h(z)=G_{p,q}^{1,p}\!\left(
 \left.\begin{matrix}a_1,\ldots,a_p\\
 b_h,b_1,\ldots,\widehat{b_h},\ldots,b_q
 \end{matrix}\right|(-1)^{p-m-n+1}z\right),
 \quad h=1,\ldots,q.
\end{equation}
The right-residue formula gives
\begin{equation}\label{eq:right-leading}
 Y_h(z)=c_hz^{b_h}(1+O(z)),\qquad c_h\ne0,
\end{equation}
under Gamma apartness.  Apply $\dd^{r}/\dd z^{r}$ for
$r=0,\ldots,q-1$.  The leading Wronskian is
\begin{equation}\label{eq:wronskian}
 \left(\prod_{h=1}^{q}c_h\right)
 z^{\sum_hb_h-q(q-1)/2}
 \prod_{1\le i<j\le q}(b_j-b_i),
\end{equation}
up to a factor $1+O(z)$.

\begin{theorem}[Corrected fundamental-system theorem]\label{thm:fundamental}
Suppose $p\le q$, the $b_h$ are pairwise distinct modulo integers, and every
leading Gamma factor in \eqref{eq:right-leading} has a quantitative apartness
witness.  Then the $q$ functions \eqref{eq:right-basis} form a fundamental
system on a sufficiently small certified punctured sector.  Dually, for
$p\ge q$, the $p$ functions obtained by placing each $a_h$ first and using the
left-residue presentation form a fundamental system when the $a_h$ are
pairwise distinct modulo integers.
\end{theorem}
\begin{proof}
The Vandermonde factor in \eqref{eq:wronskian} has a rational lower bound by
hypothesis.  Uniform tail bounds make the $O(z)$ correction smaller than half
the leading term on an explicitly searched rational radius.  Therefore the
Wronskian is apart from zero there.  Each function solves the same equation by
\cref{thm:ode}; an order-$q$ linear equation cannot have more than $q$
independent local solutions.  The usual linear-algebra proof is finite: solve
the initial-value system by Cramer's rule using the Wronskian apartness witness.
The left case follows after $z\mapsto z^{-1}$.
\end{proof}

\begin{remark}[Resonance]
If two exponents differ by an integer, the Vandermonde argument fails and the
listed $G$-functions may be dependent.  The jet construction supplies the
missing logarithmic Frobenius solutions.  Thus the source statement is correct
generically but not literally without a resonance clause.
\end{remark}

\subsection{Uniqueness without an unrestricted identity theorem}

Many later special cases are established by showing that two named functions
solve the same differential equation and have the same local coefficients.  In
this report that phrase abbreviates a finite argument: for every requested
precision choose a rational disk, compute enough coefficients, bound both
remainders, and use a finite recurrence to show coefficient equality.  If one
uses initial values instead, the Wronskian lower bound gives a finite
Gr\"onwall/Picard estimate.  No quantification over arbitrary open sets is
required.
